% =====================================================================
% Chapter 10 — Conclusions and Future Work
% =====================================================================
\chapter{Conclusions and Future Work}
\label{ch:conclusions}

This chapter concludes the dissertation. It restates what was built and what
was learned (\secref{sec:conc_summary}), maps the results back to the research
question and objectives (\secref{sec:conc_objectives}), and proposes the
experiments that would convert this project's self-consistency evidence into
externally validated accuracy (\secref{sec:conc_future}).

\section{Summary of contributions}
\label{sec:conc_summary}

This dissertation set out to answer how incomplete acoustic-sonar returns of
small underwater objects, as provided by the \suop{} dataset, can be converted
into useful 3D representations --- and to make the choice between two
reconstruction paradigms explicit. The project delivered:

\begin{enumerate}
  \item \textbf{An honest data audit.} The \suop{} release was characterised
  as a detection dataset: a fixed BlueView BV5000, an object moved by hand,
  no logged poses, truncated ping exports. Every design decision in the
  project follows from these facts (\secref{ch:dataset}).
  \item \textbf{A complete preparation pipeline.} Extraction, denoising,
  plane removal, temporal background rejection and scale-consistent clustering
  converted 120 raw cases into clean per-object segments, with a curated
  24-view subset per object for fusion (\secref{ch:dataset}).
  \item \textbf{\routeA{} --- multi-view geometric fusion.} A fully geometric
  pipeline (FPFH/RANSAC/ICP $\to$ pose graph with MST and loop closures $\to$
  fusion $\to$ Screened Poisson + BPA) ran successfully on all five
  categories: 24/24 views connected per object, fused clouds of 7--100\,k
  points, and two clean meshes per object. Two hardened variants --- a
  seabed-plane-constrained SE(2) registration and a sonar-native first-return
  ray TSDF with explicit per-vertex measurement support --- achieved median
  observed-to-surface residuals of $2.9$--$5.0\mm$ (\secref{ch:resultsA}).
  \item \textbf{\routeB{} --- single-view learned completion.} A PCN with a
  PointNet-style encoder and a coarse-to-fine folding decoder was trained with
  a symmetric Chamfer loss on sonar-realistic partials. The training data came
  from procedural models \emph{and} from \routeA{}'s own meshes, linking the
  routes. All runs converged (final training CD $1.1\times10^{-3}$--
  $2.3\times10^{-3}$), epoch sheets show category-shaped completions, and a
  real chair partial was completed at the correct scale
  (\secref{ch:resultsB}).
  \item \textbf{An explicit route rationale.} The conditions under which each
  route is the right answer, and the difference in what their outputs mean,
  were stated precisely and compared side by side (\secref{ch:comparison}).
\end{enumerate}

\section{Answers to the research question and objectives}
\label{sec:conc_objectives}

Returning to the research question of \secref{sec:research_question}:

\begin{quote}
\emph{How can incomplete acoustic-sonar returns of small underwater objects,
as provided by the \suop{} dataset, be converted into a useful 3D object
representation, and which reconstruction route should be chosen under which
data-availability conditions?}
\end{quote}

The answer, supported by the results, is twofold.

\emph{How to convert the returns:} by treating segmentation as a first-class
problem (the object must be separated from a dominant seabed), estimating
object poses from geometry rather than reading them from an encoder, and then
choosing between fusion and completion. When multiple views exist, FPFH/RANSAC/
ICP plus pose-graph optimisation reliably aligns the partials into a
self-consistent measured model. When only one view exists, a PCN trained on
sonar-realistic (partial, complete) pairs --- grounded where possible in the
geometric route's meshes --- produces a complete, category-shaped point cloud.

\emph{Which route to choose:} \routeA{} whenever overlapping views of a rigid
object are available and the output must be a faithful record of measurements
(inspection, safety assessment); \routeB{} whenever only one view exists and a
complete representation is required (classification features, simulation
targets, manipulation). The routes are complementary; their outputs differ in
completeness, in the status of unobserved regions, and in the availability of
ground truth; and they must therefore be evaluated with different protocols.

Against the six objectives of \secref{sec:research_question}: O1 (data
characterisation) is met in \secref{ch:dataset}; O2 (preparation pipeline) in
\secref{ch:dataset}; O3 (\routeA{}) in \secref{ch:routeA} and
\secref{ch:resultsA}; O4 (\routeB{}) in \secref{ch:routeB} and
\secref{ch:resultsB}; O5 (route rationale) in \secref{ch:comparison}; and O6
(evaluation and future work) across \secref{ch:resultsA}--\secref{ch:discussion}
and the remainder of this chapter.

\section{Limitations}
\label{sec:conc_limitations}

The limitations already detailed in \secref{ch:discussion} are summarised here
as the boundary of what can be claimed:

\begin{enumerate}
  \item \textbf{No external geometric ground truth.} \suop{} provides no CAD
  models, no logged object poses and no independent scan. All \routeA{}
  accuracy claims are therefore self-consistency measures, and \routeB{}
  completions cannot be verified against the real objects.
  \item \textbf{Uneven evidence.} The drum and net categories completed the
  pipelines numerically but lack dedicated rendered validation, and only
  chair/tyre received the view-aware and ray-TSDF variant treatment.
  \item \textbf{Compute-bounded exploration.} The pairwise registration cost
  and the multi-hour chair run limited the number of configurations
  (thresholds, voxels, network sizes) that could be swept.
\end{enumerate}

\section{Future work}
\label{sec:conc_future}

\subsection{A calibrated multi-view acquisition with external reference}
The single most valuable next step is to re-acquire the five \suop{} object
categories under a \emph{pose-calibrated} protocol, as proposed in the
project's acquisition specification (\texttt{SUOP\_ACQUISITION\_RECONSTRUCTION.md}):

\begin{enumerate}
  \item Keep the object rigid and fixed; move the sonar, or place the object
  on a non-reflective indexed turntable.
  \item Log the full calibrated $SE(3)$ transform for every scan (encoder or
  survey accuracy $\approx 2$~mm and $0.2^\circ$ at the 3~m working range).
  \item Acquire 24 azimuths at $15^\circ$ spacing at the 3~m range, in two
  elevation bands (the \suop{} $45^\circ$ close-range setup and a
  $15$--$25^\circ$ grazing setup), so that chair legs, seat undersides and
  tyre sidewalls are observed.
  \item Record background-only scans before and after the sequence, full
  per-ping pan/tilt/time/sound-speed metadata, and sound-speed profiles at
  start and end.
  \item Scan the same objects with a high-resolution optical or laser system
  (or use the manufacturer's CAD) to obtain an \emph{external} reference mesh.
\end{enumerate}

With logged poses, the pose-estimation stage of \routeA{} is bypassed and the
raw first-return rays can be integrated directly into a 3--5~mm TSDF, giving a
\emph{measured} reference. With the external scan, both routes can be
evaluated quantitatively on the same objects:

\begin{itemize}
  \item \routeA{}: Chamfer distance and F-score between the fused/TSDF surface
  and the external mesh; registration error against the logged poses; coverage
  as the fraction of the external surface within a tolerance.
  \item \routeB{}: Chamfer distance and F-score of the completions against the
  external mesh on \emph{held-out real partials} --- this is the definitive
  test of whether the learned prior generalises to genuine sonar rather than
  simulated partials.
\end{itemize}

\subsection{Hardening each route}
Independent of the new acquisition, each route can be hardened:

\begin{itemize}
  \item \routeA{}: integrate the later-return (multipath) evidence through
  cross-view consistency instead of discarding it at source; add an explicit
  outlier-rejection pass in the TSDF; and expose per-vertex coverage as a
  standard output (already prototyped in the ray-TSDF variant).
  \item \routeB{}: evaluate against held-out \emph{real} partials with
  F-score; replace the fixed 4096-point fine output with a
  density-adaptive decoder (e.g.\ FoldingNet-style
  \citep{yang2018folding} or a transformer-based completion
  \citep{xie2020pointr}); and train a multi-category model that reports
  category uncertainty.
\end{itemize}

\subsection{Chaining the routes}
Finally, the mesh-conditioned training result suggests a promising hybrid:
use \routeA{} to build a library of measured meshes per category, use those
meshes (not procedural models) to train \routeB{}, and use \routeB{}'s
completions as \emph{priors} inside a revised \routeA{} (e.g.\ as soft
regularisers in the TSDF). This would give the geometric route the
completeness it lacks and the learned route the sensor-grounded prior it
needs --- the most productive direction this dissertation can point towards.

\section{Closing remark}
\label{sec:conc_closing}

Underwater sonar is a hard sensor: its point clouds are partial, noisy and
unevenly sampled, and the \suop{} dataset adds the further difficulty of
unlogged object pose. This dissertation shows that two complementary routes
can nevertheless convert such returns into useful 3D representations --- one
that faithfully fuses what was measured when several views exist, and one that
inferentially completes what was not observed when only a single view exists.
The two routes are not competitors; they are the two halves of a complete
answer, and the next acquisition campaign with logged poses and external
reference will let both halves be measured against the truth they ultimately
serve.
